% TOP_PROBLEM: {"id": 30007024, "title": "Conditional abelianity for special half-weight fourfold pairs", "category_id": 5, "category": {"id": 5, "name": "algebraic_geometry", "display_name": "Algebraic Geometry", "description": "Geometric objects defined by polynomial equations.", "slug": "algebraic-geometry", "order_index": 5, "created_at": "2026-07-31T15:26:25.671Z"}, "status": "open"}
% FIRST_POSTED: 2026-10-06
% ATTEMPT_STATUS: solved
% OPENAI_CLAIM: {"schema_version":1,"source_repo":"https://github.com/openai/math","source_commit":"adc7f1241b42e322a6451854ab7e4b4c146bf78a","release_date":"2026-10-06","match":"full","resolution":"proved","independently_reviewed":false,"second_pass":"agreed","families":[{"family":"057","title":"Fundamental groups of special complex varieties and root orbifolds","subject":"Algebraic and complex geometry","manuscripts":[{"title":"A conditional abelianity theorem for special fourfold pairs with a half-weight divisor","pdf_path":"preprints/A-conditional-abelianity-theorem-for-special-fourfold-pairs-with-a-half-weight-divisor-October-5-2026/main.pdf","readme_path":"preprints/A-conditional-abelianity-theorem-for-special-fourfold-pairs-with-a-half-weight-divisor-October-5-2026/README.md","date":"2026-10-05","theorem_ref":"thm:main (source/sections/introduction.tex:39)","statement_sources":[{"path":"preprints/A-conditional-abelianity-theorem-for-special-fourfold-pairs-with-a-half-weight-divisor-October-5-2026/build/source/sections/introduction.tex","line":39,"label":"thm:main"}]}],"lean":null,"reasoning_trace":null}]}
\documentclass[11pt]{article}
\usepackage[margin=25mm]{geometry}
\usepackage{fontspec}
\setmainfont{Latin Modern Roman}
\usepackage{amsmath,amssymb,mathtools}
\usepackage{unicode-math}
\setmathfont{Latin Modern Math}
\usepackage{enumitem}
\usepackage{xurl,hyperref,bookmark}
\hypersetup{colorlinks=true,urlcolor=blue}
\setcounter{secnumdepth}{0}
\setlength{\parindent}{0pt}
\setlength{\parskip}{5pt}
\setlength{\emergencystretch}{3em}
\begin{document}
\section*{Conditional abelianity for special half-weight fourfold pairs}

\subsection{OpenAI's claim}
Subject to the stated companion abelianity input, virtual abelianity is claimed for the full fundamental group of a special order-two root orbifold from a smooth projective fourfold and one smooth connected divisor. The cited manuscript argues for that result; its relation to this problem is specified below.

\subsection{Mathematical statement of the claimed result}
\paragraph{Mathematical claim 1.} Assume every smooth connected compact Kähler special manifold has virtually abelian fundamental group in every dimension. Special means there is no line bundle $L$ with a nonzero $L\to\Omega^p$ and $\kappa(L)=p$. Let $X$ be a smooth connected projective complex fourfold and $D$ a reduced divisor that is smooth, nonempty and connected. Require the order-two root orbifold $\sqrt[2]{(X,D)}$ to be special by the same equivariant differential criterion; locally its chart has $z_1=w_1^2$ and sign stabilizer. OpenAI claims
\[
\pi_1(X\setminus D,x)/\langle\!\langle\gamma^2\rangle\!\rangle\text{ is virtually abelian},
\]
where $\gamma$ is a positively oriented transported meridian. The ordinary-manifold abelianity theorem is an explicit premise; a general orbifold abelianity theorem is not assumed or concluded.

\subsection{Outline of the argument}
Manuscript 1 (family 057, 2026-10-05): A smooth projective eightfold over the root orbifold has genus-one product fibers and submersive lifts over every base divisor. Those lifts detect poles of descended differential tensors, so specialness passes to the source, while even pullback multiplicities give a surjection of its fundamental group onto the orbifold group.

\subsection{Scope relative to this problem}
Manuscript 1 (family 057, 2026-10-05): The problem selects the mathematical statement summarized here, including its stated hypotheses and parameter range. The cited theorem gives a claimed proof of that assertion in full.

\subsection{Formal verification}
Family 057 has no Lean scope page in this snapshot; no formalization of these manuscript results is listed.

This site has not executed the Lean files or checked their dependencies. These descriptions report the scope of the linked declarations, rather than a new verification of the proofs.

\subsection{Status caveat}
OpenAI attributes these results to an internal model. Its release notice cautions that results lacking full formalization can contain mistakes. This site has not independently assessed the proofs; a solved label records the scope of the reported claim.

\subsection{Sources}
\begin{itemize}
\item OpenAI, \emph{A conditional abelianity theorem for special fourfold pairs with a half-weight divisor}, 2026-10-05; citation key \texttt{OAI:A-conditional-abelianity-theorem-for-special-fourfold-pairs-with-a-half-weight-divisor-October-5-2026}. \href{https://github.com/openai/math/blob/adc7f1241b42e322a6451854ab7e4b4c146bf78a/preprints/A-conditional-abelianity-theorem-for-special-fourfold-pairs-with-a-half-weight-divisor-October-5-2026/main.pdf}{PDF (GitHub)}; \href{https://github.com/openai/math/raw/adc7f1241b42e322a6451854ab7e4b4c146bf78a/preprints/A-conditional-abelianity-theorem-for-special-fourfold-pairs-with-a-half-weight-divisor-October-5-2026/main.pdf}{PDF (direct)}; \href{https://github.com/openai/math/blob/adc7f1241b42e322a6451854ab7e4b4c146bf78a/preprints/A-conditional-abelianity-theorem-for-special-fourfold-pairs-with-a-half-weight-divisor-October-5-2026/README.md}{Manuscript page and citation}; \href{https://github.com/openai/math/blob/main/preprints/A-conditional-abelianity-theorem-for-special-fourfold-pairs-with-a-half-weight-divisor-October-5-2026/main.pdf}{Latest version}; Statement source: \href{https://github.com/openai/math/blob/adc7f1241b42e322a6451854ab7e4b4c146bf78a/preprints/A-conditional-abelianity-theorem-for-special-fourfold-pairs-with-a-half-weight-divisor-October-5-2026/build/source/sections/introduction.tex\#L39}{thm:main}.
\item \href{https://github.com/openai/math/blob/adc7f1241b42e322a6451854ab7e4b4c146bf78a/CONTENTS.md}{OpenAI catalogue entry}.
\end{itemize}
\end{document}
